\section{Intrinsic dynamics}\label{app:intrinsic}
Special conserved spectral values express some covariance blocks in terms
of fewer variables. This appendix derives these exact reductions.

\subsection{A two-point signed spectrum: hyperbolic isotropy}
A special two-point spectrum determines the hidden covariance blocks
from the predictor and therefore closes a smaller equation.
\begin{definition}[Hyperbolic isotropy]\label{def:hyperbolic-isotropy}
For \(\alpha>0\), the hyperbolically isotropic (HI) leaf is
\begin{equation}\label{eq:intrinsic-hi}
 \mathcal M_\alpha:=\{\Sigma\succeq0:\Sigma J\Sigma=\alpha^2J\}.
\end{equation}
\end{definition}
The identity forces $\Sigma\succ0$ and is equivalent to
$(J\Sigma)^2=\alpha^2\Id_m$. The symmetric representative of $J\Sigma$
then has $d_u$ eigenvalues $+\alpha$ and $d_v$ eigenvalues $-\alpha$;
equivalently, $\mathcal H_k=\alpha^k(d_u+(-1)^kd_v)$ for $k=1,\ldots,m$.
Spectral conservation proves invariance of $\mathcal M_\alpha$ despite
the repeated eigenvalues. Directly, the derivative of $\Sigma J\Sigma$
is $-G(\Sigma J\Sigma)-(\Sigma J\Sigma)G$, with constant solution $\alpha^2J$.

\begin{proposition}[Predictor-only intrinsic equation]\label{prop:intrinsic-hi}
Every predictor \(\Svu\in\R^{d_v\times d_u}\) has a unique covariance
on \(\mathcal M_\alpha\), given by
\begin{equation}\label{eq:intrinsic-graph}
 \Sigma=
 \begin{pmatrix}
 ((\Svu)^\top\Svu+\alpha^2\Id_{d_u})^{1/2}&(\Svu)^\top\\
 \Svu&(\Svu(\Svu)^\top+\alpha^2\Id_{d_v})^{1/2}
 \end{pmatrix}.
\end{equation}
On this leaf, the Bures dynamics close on the predictor:
\begin{equation}\label{eq:intrinsic-hyperbolic-dynamics}
 \dSvu=-E((\Svu)^\top\Svu+\alpha^2\Id)^{1/2}
       -(\Svu(\Svu)^\top+\alpha^2\Id)^{1/2}E,
 \qquad E:=\nabla\ell(\Svu).
\end{equation}
For \(d_u=d_v=1\), this is
\(\dSvu=-2\sqrt{(\Svu)^2+\alpha^2}\,\ell'(\Svu)\).
\end{proposition}
\begin{proof}
Expanding \eqref{eq:intrinsic-hi} gives
\[
 (\Suu)^2-(\Svu)^\top\Svu=\alpha^2\Id,\quad
 (\Svv)^2-\Svu(\Svu)^\top=\alpha^2\Id,\quad
 \Svu\Suu=\Svv\Svu.
\]
Positivity selects the positive square roots. Conversely, an SVD proves
the intertwining identity and gives eigenvalues \(\alpha\) and
\(\sqrt{\sigma_i^2+\alpha^2}\pm\sigma_i>0\), proving membership in
\(\mathcal M_\alpha\). Substituting the diagonal blocks into
\eqref{eq:block-flow-general} gives the intrinsic equation.
\end{proof}
The same eigenvalue calculation gives the exact extremal envelopes
\[
 \lambda_{\max}(\Sigma)=\sqrt{\norm{\Svu}_{\op}^2+\alpha^2}
                         +\norm{\Svu}_{\op},\qquad
 \lambda_{\min}(\Sigma)=\frac{\alpha^2}{\lambda_{\max}(\Sigma)}.
\]
In the general convergence theorems this special spectral level has
\(\aplus{\Sigma_0}=\aminus{\Sigma_0}=\alpha\). Thus its loss exponent is \(4\alpha\kappa\),
and the GD contraction factor is \(1-\sqrt2\alpha\kappa\tau\).

\paragraph{Finite-width form.}
The covariance constraint has an exact finite-width form that depends
only on the product $WW^\top$.
For \(\Sigma=WW^\top/n\), condition \eqref{eq:intrinsic-hi} is
\(WW^\top JWW^\top=n^2\alpha^2J\). In terms of \(U,V\), it reads
\[
 \begin{gathered}
 (UU^\top)^2-UV^\top VU^\top=n^2\alpha^2\Id_{d_u},\qquad
 (VV^\top)^2-VU^\top UV^\top=n^2\alpha^2\Id_{d_v},\\
 VU^\top UU^\top=VV^\top VU^\top.
 \end{gathered}
\]
Any covariance in \eqref{eq:intrinsic-graph} admits empirical factors at
every width \(n\ge m\), for example \(W=\sqrt n\,\Sigma^{1/2}Q_0\)
with \(Q_0Q_0^\top=\Id_m\).

\paragraph{Comparison with relaxed balance.}
\citet[Theorem 3.8]{MarcottePeyreGribonval2026} derive intrinsic dynamics
under a different condition,
\(\Gamma_{\rm rb}:=(V^\top V-U^\top U)/n=\beta\Id_n\),
on a full-column-rank stratum. It implies
\[
 (\Suu)^2+\beta\Suu=(\Svu)^\top\Svu,\qquad
 (\Svv)^2-\beta\Svv=\Svu(\Svu)^\top.
\]
For \(\beta\ne0\), zero singular directions have ambiguous branches.
Where both factors have full column rank, the range projectors
\(\Pi_u,\Pi_v\) of \((\Svu)^\top,\Svu\) select
\[
 \begin{aligned}
 \Suu&=\Pi_u[-\tfrac\beta2\Id+
             ((\Svu)^\top\Svu+\tfrac{\beta^2}{4}\Id)^{1/2}],\\
 \Svv&=\Pi_v[\tfrac\beta2\Id+
             (\Svu(\Svu)^\top+\tfrac{\beta^2}{4}\Id)^{1/2}].
 \end{aligned}
\]
If \(n=d_u=d_v\) and both factors are invertible, the projectors are
identities and the predictor equation coincides with
\eqref{eq:intrinsic-hyperbolic-dynamics} for \(\alpha=|\beta|/2\).
The leaves remain distinct: for \(\beta\ne0\), relaxed balance is equivalent to
\(\Sigma J\Sigma=-\beta\Sigma\), \(\rankop\Sigma=n\).
Indeed a left inverse of \(W/\sqrt n\) proves the converse.
It forces \(n\le d_v\) if \(\beta>0\), and \(n\le d_u\) if \(\beta<0\);
positive HI instead has rank \(m\).
For \(\beta=0\), positivity selects the unshifted square roots without a
rank assumption, giving the boundary equation \(\alpha=0\), not a
positive-definite initialization.
Without the rank qualifications, relaxed balance does not ensure
predictor closure: \(n=1,d_u=d_v=2,U=0,V=e_1\) or \(e_2\) gives the
same predictor and \(\beta=1\), but different predictor velocities for
\(\ell(S):=\norm{S+\Id_2}_{\F}^2/2\).

\begin{remark}[Conserved values do not generally close the predictor]
For \(d_u=d_v=2\), take \(\Svu=0\), \(\Svv=\diag(3,4)\), and
\(\Suu=\diag(1,2)\) or \(\diag(2,1)\).
Both states have signed spectrum \(\{1,2,-3,-4\}\), hence the same
spectral laws. For \(\ell(S):=\norm{S+e_1e_1^\top}_{\F}^2/2\),
their predictor velocities are \(-4e_1e_1^\top\) and
\(-5e_1e_1^\top\).
The special two-point spectrum gives an algebraic graph unavailable
at these general spectral values.
\end{remark}

\subsection{A single phase spectral scale: conformal symplectic isotropy}
A special conformal spectral identity eliminates the velocity covariance
within a positive-position chart of the phase flow.
\begin{definition}[Conformal symplectic isotropy]\label{def:csi}
A phase covariance is conformally symplectically isotropic (CSI) at
scale \(\theta\ge0\) if
\begin{equation}\label{eq:intrinsic-csi}
 \Phase \Jcan\Phase=\theta^2\Jcan.
\end{equation}
\end{definition}
For $\theta>0$, this forces $\Phase\succ0$ and is equivalent to
$\Phase^{1/2}\Jcan\Phase^{1/2}$ having $m$ copies of each eigenvalue
$+i\theta,-i\theta$.
The renormalized trace powers in Proposition~\ref{prop:phase-laws}
are zero at odd orders and $2m(-1)^k\theta_0^{2k}$ at order $2k$.
Conformal transport preserves \eqref{eq:intrinsic-csi} with
$\theta_t=e^{-\gamma t}\theta_0$, including the singular boundary
$\theta_0=0$, which is specified algebraically, without diagonalizability.

\begin{proposition}[Reduced phase dynamics]\label{prop:intrinsic-csi}
On the chart \(A:=\Spos\succ0\), condition \eqref{eq:intrinsic-csi}
is equivalent to
\[
 D=C=C^\top,\qquad Q=\theta^2A^{-1},
 \qquad \Xi=AC,\quad \Svel=CAC+\theta^2A^{-1}.
\]
The reduced equations are
\begin{equation}\label{eq:csi-reduced}
 \dot A=CA+AC,\qquad
 \dot C=-G-\gamma C-C^2+\theta_t^2A^{-2},
 \qquad\dot\theta_t=-\gamma\theta_t.
\end{equation}
They use only the two symmetric matrices \(A,C\), with known scale
\(\theta_t\), instead of all blocks of the phase covariance.
\end{proposition}
\begin{proof}
Substitute \(\Xi=AD^\top\), \(\Svel=DAD^\top+Q\).
The upper-left block of \eqref{eq:intrinsic-csi} is
\(A(D-D^\top)A=0\), forcing symmetry.
The upper-right block then becomes \(AQ=\theta^2\Id\);
the other blocks follow.
Substitution in Proposition~\ref{prop:regression} gives
\eqref{eq:csi-reduced}. This calculation includes \(\theta=0\).
\end{proof}
White position and velocity blocks with \(\Xi_0=0\) give
\(\theta_0=\sigma_{vu}\sigma_{qp}\). Zero velocity gives
\(C_0=0,\theta_0=0\), allowing any \(A_0\succ0\).
Neither choice is needed for convergence: the general theorem controls
\(D\) by energy even when it is nonsymmetric and \(Q\) is independent.

\begin{remark}[What the reduced state still contains]
Even on CSI, the two cross blocks \(\Svu,\Sqp\) do not generally form a
closed state. With \(\theta=1,C=0\), choose \(A=\Id_2\) or \(2\Id_2\)
and \(\Svel=A^{-1}\). Both have
\(\Svu=\Sqp=\dSvu=0\), yet
\(\ddSvu=-2E\) or \(-4E\) for the same nonzero scalar loss gradient.
Thus the full position block remains part of the intrinsic variables.
\end{remark}
